\begin{proof}
\paragraph{1. Degenerate cases.}
If \(a=0\) then the right–hand side equals \(\dfrac{b^{q}}{q}\ge 0\) while the left–hand side is \(0\); the inequality holds.  
If \(b=0\) the argument is symmetric.  
Hence we may assume \(a>0\) and \(b>0\), so that all logarithms below are well defined.

\paragraph{2. Weighted AM–GM via Jensen’s inequality.}
Consider the convex function \(f(t)=e^{t}\) on \(\mathbb{R}\) (since \(f''(t)=e^{t}>0\)).  
Apply Jensen’s inequality to the two points  

\[
x_{1}= \ln a^{p},\qquad x_{2}= \ln b^{q},
\]

with weights  

\[
\lambda_{1}= \frac{1}{p},\qquad \lambda_{2}= \frac{1}{q}.
\]

Because \(\lambda_{1},\lambda_{2}\ge 0\) and \(\lambda_{1}+\lambda_{2}= \frac1p+\frac1q =1\), Jensen gives  

\[
f\!\Bigl(\lambda_{1}x_{1}+\lambda_{2}x_{2}\Bigr)
\le
\lambda_{1}f(x_{1})+\lambda_{2}f(x_{2}).
\tag{1}
\]

Writing the weights explicitly we obtain  

\begin{align*}
e^{\frac1p\ln a^{p}+\frac1q\ln b^{q}}
&\le
\frac1p\,e^{\ln a^{p}}+\frac1q\,e^{\ln b^{q}} .
\end{align*}
\tag{2}

\paragraph{3. Simplification.}
For any positive \(x\) we have \(e^{\ln x}=x\). Moreover  

\[
e^{\frac1p\ln a^{p}} = a,\qquad
e^{\frac1q\ln b^{q}} = b,
\]

since \(\frac1p\ln a^{p}= \ln a\) and \(\frac1q\ln b^{q}= \ln b\).  
Substituting these equalities into (2) yields  

\[
ab \;\le\; \frac{a^{p}}{p}+\frac{b^{q}}{q}.
\tag{3}
\]

Thus the desired inequality holds for all \(a,b>0\). Together with the trivial cases from Step 1, (3) is valid for all non‑negative \(a,b\).

\paragraph{4. Equality condition.}
For a strictly convex function, Jensen’s inequality is an equality **iff** the
arguments are equal. Hence equality in (1) (and therefore in (3)) occurs exactly when  

\[
x_{1}=x_{2}\quad\Longleftrightarrow\quad
\ln a^{p}= \ln b^{q}\quad\Longleftrightarrow\quad
a^{p}=b^{q}.
\]

When \(a=b=0\) both sides of (3) are zero, so equality also holds. No other
choice of \((a,b)\) gives equality. Consequently the equality case is precisely  

\[
a^{p}=b^{q}\qquad(\text{including }a=b=0).
\]

\end{proof}